% !TeX root = ../../Book3.tex
\chapter*{后记}
\addcontentsline{toc}{chapter}{后记}

本卷从复导数开始，最后到达素数的渐近分布。两端之间并非一条只向前推进的直线：Cauchy 积分在留数、边值与延拓中反复出现，零点计数经过 Jensen 公式又进入整函数和模形式，紧性则从正规族延伸到共形映射和覆盖。每次重逢，所保存的局部信息相似，面对的整体障碍却不同。

\ProseParagraphBreak

写作中最需要保持的区别，是局部计算与全局存在之间的区别。一个可积的微分形式未必有全局单值原函数，局部逆未必拼成整体逆，边界上的点值也未必能由区域内部无条件恢复。拓扑与边界假设因此没有被藏进图形直觉。读者若能在使用公式以前先说明路径、分支、区域与收敛方式，许多看似意外的矛盾就会消失。

\ProseParagraphBreak

本卷也有意把几个深层结论的边界写出来。统一化、Riemann--Roch、可测 Riemann 映射等理论各自需要相当长的准备，不能由几行说明代替。相应位置给出具体版本、用途和证明范围，使读者能够继续查阅，也能辨认后续哪一步依赖这一输入。其余能够在当前工具下完成的推导，则不因最终目标深远而省略基础环节。

\ProseParagraphBreak

读完之后，不妨换一种次序再看几段证明。先从一个多连通区域出发，比较原函数、调和共轭与覆盖各自遇到的障碍；再把 Jensen 公式用于一个显式整函数，检查零点密度怎样限制增长；最后从 Dirichlet 级数的收敛半平面出发，追踪解析延拓、边界零点和 Tauber 型论证怎样共同产生算术渐近。这样回读，较容易发现自己掌握的是公式，还是公式背后的条件。

\ProseParagraphBreak

继续进入专著时，共形几何会引向曲面与拟共形变形，势论会引向容量和精细边界行为，模形式会引向更多群上的自守对象与 \(L\)-函数。它们的语言不断扩大，但“局部信息怎样限制整体”仍是值得保留的问题。遇到一个强定理时，先找最简单的模型和失败的边界，往往比急着追求最一般形式更有收获。

\ProseParagraphBreak

三卷至此把测度、频率与复结构放到了同一段学习经历中。它们各有自己的方法，也共享对极限、紧性、表示和假设的认真态度。书页可以到这里结束，读者自己的例子、反例和证明却应从这里继续增加。

\bigskip
\begin{flushright}
R. EverStarine

2026年10月11日
\end{flushright}
